\newpage

\section{Autoregressive Models}

% ============================================================
\subsection*{Chain Rule of Probability: How and Why It Works}
% ============================================================

\subsubsection*{1. Where it comes from: conditional probability}

Everything starts from the definition of conditional probability:
\[
p(A \mid B) = \frac{p(A, B)}{p(B)}, \quad \text{for } p(B) > 0
\]
Read it as: ``Among all outcomes where $B$ happened, what fraction also has $A$?''

Multiply both sides by $p(B)$:
\[
p(A, B) = p(B)\, p(A \mid B)
\]
This is the \textbf{two-variable chain rule}. It is just the definition rearranged, not a new law.

\subsubsection*{2. Extending to many variables}

Apply the same step repeatedly, treating everything before the current variable as ``$B$''.

For three variables, group $(x_1, x_2)$ as one block:
\[
p(x_1, x_2, x_3) = p(x_1, x_2)\, p(x_3 \mid x_1, x_2)
\]
Then split $p(x_1, x_2)$ again:
\[
p(x_1, x_2, x_3) = p(x_1)\, p(x_2 \mid x_1)\, p(x_3 \mid x_1, x_2)
\]
Repeating this $D$ times gives the general form:
\[
p(x_1, \dots, x_D) = \prod_{i=1}^{D} p(x_i \mid x_{<i})
\]

\subsubsection*{3. Why it works: the product telescopes}

Write each conditional using its definition:
\[
p(x_1) \cdot \frac{p(x_1, x_2)}{p(x_1)} \cdot \frac{p(x_1, x_2, x_3)}{p(x_1, x_2)}
\]
Each numerator cancels the next denominator. Only $p(x_1, x_2, x_3)$ survives.

So the chain rule is an \textbf{identity}. It holds for every distribution and makes no assumptions:
\begin{itemize}
  \item \textbf{No independence is assumed.} Each variable may depend on all earlier ones.
  \item \textbf{Nothing is approximated.} The product equals the joint exactly.
\end{itemize}
The only requirement is that each prefix has nonzero probability, so the divisions are defined.

\subsubsection*{4. Intuition: narrowing down step by step}

Think of the joint probability as a fraction of all possible outcomes. You reach it by successive filtering.

Take a language model generating the sentence \textbf{``the cat sat''}:

\begin{center}
\begin{tabular}{c l p{8cm} c}
\toprule
Step & Factor & Meaning & Value \\
\midrule
1 & $p(\text{the})$ & Fraction of all sentences that start with ``the'' & 0.05 \\
2 & $p(\text{cat} \mid \text{the})$ & Of those, fraction whose next token is ``cat'' & 0.01 \\
3 & $p(\text{sat} \mid \text{the, cat})$ & Of those, fraction whose next token is ``sat'' & 0.20 \\
\bottomrule
\end{tabular}
\end{center}

\[
p(\text{the cat sat}) = 0.05 \times 0.01 \times 0.20 = 0.0001
\]
Each factor is a ``fraction of a fraction.'' Multiplying them gives the share of the whole sentence space occupied by exactly this sequence.

The same logic holds for images. The probability of a whole image equals:
\begin{itemize}
  \item the probability of the first pixel's intensity,
  \item times the probability of the second pixel given the first,
  \item and so on, up to the last pixel given all pixels before it.
\end{itemize}

\subsubsection*{5. Contrast: what would be wrong}

If you assumed the tokens were independent, you would compute:
\[
p(\text{the})\, p(\text{cat})\, p(\text{sat})
\]
This ignores the fact that ``sat'' is far more likely after ``the cat'' than in general. The joint would be badly wrong.

The chain rule keeps every dependency by conditioning on the full history.

\subsubsection*{6. Order does not matter for correctness}

The derivation never used a particular order. You could factor right-to-left, or in any permutation:
\[
p(x_1, x_2, x_3) = p(x_3)\, p(x_2 \mid x_3)\, p(x_1 \mid x_2, x_3)
\]
All orderings give the \textbf{same joint}. What changes is the set of conditionals. Some orderings are easier for a network to learn, such as left-to-right for text, time order for audio, and raster order for images. That is why the lecture says ``the order is part of the model.''

\subsubsection*{7. Bridge to autoregressive models}

The chain rule turns one huge, intractable object into $D$ small conditionals. Examples of that huge object are a distribution over all $32 \times 32 \times 3$ images, or over all possible sentences.

An autoregressive model's only job is to learn those conditionals with a shared network. Because the factorization is exact, the model's likelihood is exact too.

% ============================================================
\subsection*{Meaning of ``The Chain Rule Is an Identity''}
% ============================================================

An \textbf{identity} is an equation that is true for every possible value of its variables, not just under certain conditions. For example, $(a+b)^2 = a^2 + 2ab + b^2$ holds for all numbers $a, b$; you never have to check whether it applies.

``The chain rule is an identity'' means:
\[
p(x_1, \dots, x_D) = \prod_{i=1}^{D} p(x_i \mid x_{<i})
\]
is true for \textbf{every} probability distribution and \textbf{every} sequence. It follows directly from the definition of conditional probability, as the telescoping cancellation showed, so it cannot fail.

Contrast this with a \textbf{modeling assumption}, which can be wrong:
\begin{itemize}
  \item \textbf{Identity (always true):} $p(\text{the cat sat}) = p(\text{the})\,p(\text{cat} \mid \text{the})\,p(\text{sat} \mid \text{the, cat})$
  \item \textbf{Assumption (may be false):} $p(\text{the cat sat}) = p(\text{the})\,p(\text{cat})\,p(\text{sat})$, which claims the tokens are independent
\end{itemize}

\textbf{Why it matters for autoregressive models:} writing the joint as a product of conditionals loses no generality. Any distribution over images, text, or audio can be represented this way.

The only possible source of error is the network's estimate of each conditional $p_\theta(x_i \mid x_{<i})$, not the factorization itself. That is why the lecture calls the likelihood ``exact'': there is no approximation gap like the ELBO in VAEs.

% ============================================================
\subsection*{Chain Rule for Dependent vs.\ Independent Variables}
% ============================================================

The chain rule does \textbf{not} assume dependence. It \textit{allows} dependence but does not \textit{require} it. It holds in both cases.

\subsubsection*{What happens when variables are independent}

Independence means that knowing the past tells you nothing about the next variable:
\[
p(x_i \mid x_{<i}) = p(x_i)
\]
Substitute this into the chain rule:
\[
p(x_1, \dots, x_D) = \prod_{i} p(x_i \mid x_{<i}) = \prod_{i} p(x_i)
\]
This is exactly the familiar product rule for independent variables. So the independence formula is just a \textbf{special case} of the chain rule, where each conditional simplifies to a marginal.

\subsubsection*{Generative example}

Suppose a model generates \textbf{white-noise audio}, where each sample is drawn independently from $\mathcal{N}(0, \sigma^2)$:
\begin{itemize}
  \item The chain rule still gives $p(x_1)\,p(x_2 \mid x_1)\,p(x_3 \mid x_1, x_2)\cdots$
  \item Because the samples are independent, each $p(x_i \mid x_{<i})$ equals $\mathcal{N}(0, \sigma^2)$.
  \item The result is a product of identical Gaussians, which is correct.
\end{itemize}

Now suppose it generates \textbf{speech}, where each sample depends strongly on the previous ones:
\begin{itemize}
  \item The conditionals carry that dependence.
  \item The chain rule still holds, and now the conditioning actually matters.
\end{itemize}

\subsubsection*{The one real condition}

The only requirement is that each prefix must have nonzero probability, so $p(x_i \mid x_{<i})$ is defined.

This condition is about the conditional being well-defined. It has nothing to do with dependence or independence. For a valid sequence you are evaluating, it is satisfied automatically.

\subsubsection*{Summary}

\begin{center}
\begin{tabular}{l c l}
\toprule
Situation & Chain rule holds? & What the conditionals look like \\
\midrule
Dependent (text, speech, images) & Yes & Depend on the history \\
Independent (white noise) & Yes & Reduce to plain marginals \\
\bottomrule
\end{tabular}
\end{center}

This is why it is called an identity: it is true regardless of the dependence structure. Dependence only changes what the conditionals turn out to be.

% ============================================================
\subsection*{Definition of an Autoregressive Model}
% ============================================================

From the lecture: \textit{``An autoregressive (AR) model represents the joint by the chain-rule product and parameterizes each conditional with a shared function of the past.''}

\subsubsection*{The sentence, broken into parts}

\textbf{1. ``represents the joint by the chain-rule product''}

The model's goal is the joint distribution $p(\mathbf{x})$, for example a distribution over all possible images or sentences. Instead of modeling it directly, the model writes it as
\[
p(\mathbf{x}) = \prod_{i=1}^{D} p(x_i \mid x_{<i})
\]
Since the chain rule is an identity, this loses nothing. The problem has changed from ``model the whole object at once'' to ``model $D$ next-step predictions.''

\textbf{2. ``parameterizes each conditional''}

Each factor $p(x_i \mid x_{<i})$ is unknown, so a neural network with parameters $\theta$ is used to produce it. ``Parameterize'' means the conditional's form is set by learnable numbers $\theta$, written $p_\theta(x_i \mid x_{<i})$.

\textbf{3. ``with a shared function''}

\textbf{One} network with \textbf{one} set of parameters $\theta$ produces every conditional. You do not train a separate network for position 1, position 2, and so on.

Sharing is necessary, not just convenient. Take a $32 \times 32 \times 3$ image with 256 intensity values:
\begin{itemize}
  \item Without sharing, the conditional for position $i$ would need an entry for every possible prefix, which is $256^{i-1}$ cases.
  \item At position 3072, that number is astronomically large, so no lookup table could store it.
\end{itemize}
A shared network avoids this. It learns general rules, such as ``pixels near an edge tend to continue that edge,'' and applies them at every position.

\textbf{4. ``of the past''}

The network may look \textbf{only} at $x_1, \dots, x_{i-1}$. It must never see $x_i$ itself or anything after it. This is the causal constraint. If the network saw $x_i$, it could simply copy the answer, and the product would no longer equal the true joint.

\subsubsection*{The equation}

\[
p_\theta(x_i \mid x_{<i}) = f_\theta(x_1, \dots, x_{i-1})
\]

\begin{center}
\begin{tabular}{l p{11cm}}
\toprule
Symbol & Meaning \\
\midrule
$f_\theta$ & The neural network: an RNN, masked CNN, or Transformer \\
$\theta$ & Its parameters, the same at every position $i$ \\
Input $x_1, \dots, x_{i-1}$ & The past only, which is the context \\
Output & A \textbf{probability distribution} over the possible values of $x_i$ \\
\bottomrule
\end{tabular}
\end{center}

The key point is that $f_\theta$ outputs a \textbf{distribution}, not a single predicted value.

\textbf{Example: generating text.} Suppose the context is ``the cat'' and the vocabulary has 50{,}000 tokens. Then $f_\theta(\text{the, cat})$ outputs a 50{,}000-entry probability vector via softmax:
\begin{itemize}
  \item ``sat'' $\rightarrow$ 0.20
  \item ``is'' $\rightarrow$ 0.15
  \item ``ran'' $\rightarrow$ 0.08
  \item \dots and so on, summing to 1.
\end{itemize}
Uses of this vector:
\begin{itemize}
  \item \textbf{Training:} read off the probability of the true next token and add its $-\log$ to the loss.
  \item \textbf{Generation:} sample a token from this vector.
\end{itemize}

\textbf{Example: generating an image pixel.} The context is all pixels above, plus those to the left in the current row. $f_\theta$ outputs a 256-way softmax over the intensities of the next channel value.

For an audio sample, it could instead output the parameters of a Gaussian or a mixture of logistics.

\subsubsection*{The ``auto'' part}

``Autoregressive'' means the variable is predicted (\textit{regressed}) from earlier values of \textbf{itself} (\textit{auto}). Pixels predict pixels, tokens predict tokens, audio samples predict audio samples. There is no separate input stream.

This differs from, say, a translation model, which predicts French from an English input. The conditional PixelCNN in the lecture adds an external $\mathbf{h}$, but the autoregressive core is still conditioning on the data's own past.